%! AP Statistics — Unit 2, Lesson 2.3 — Student Guided Notes
\documentclass[10pt]{article}
\usepackage{apstats-article}
\usepackage{apstats-boxes}

%=============================================================================
\begin{document}

\pageheader{Unit 2, Lesson 2.3}{Guided Notes}
\namedateperiod

\vspace{0.12in}

% ── OBJECTIVE ────────────────────────────────────────────────────────────────
\begin{objectivebox}
\small By the end of this lesson, I will be able to\dots\\[4pt]
\writeline\\[1pt]
\writeline\\[1pt]
\writeline
\end{objectivebox}

\vspace{0.1in}

% ── VOCABULARY ───────────────────────────────────────────────────────────────
\begin{vocabbox}
\termblanklong{Conditional proportion ($\hat{p}$)}
\termblanklong{Difference in proportions}
\termblanklong{Relative risk (risk ratio)}
\termblanklong{Baseline group}
\end{vocabbox}

\vspace{0.1in}

% ── HOOK ─────────────────────────────────────────────────────────────────────
\begin{hookbox}
\small\textbf{Scenario:} A clinical trial tested a new vaccine. Of 400 vaccinated participants,
32 contracted the flu. Of 400 unvaccinated participants, 80 contracted the flu.

\vspace{6pt}
\textbf{Statistical question:} \textit{How much does vaccination reduce the likelihood of getting the flu?}

\vspace{6pt}
\textbf{Discussion:}
\begin{enumerate}[leftmargin=1.5em, itemsep=2pt]
  \item We already know the conditional proportions from Lesson 2.2. What are they here?\\
        Vaccinated: $\hat{p}_V = \blank{2cm}$ \quad Unvaccinated: $\hat{p}_U = \blank{2cm}$
  \item The proportions differ --- but \emph{how} do we describe that difference with a single number?\\
        \writeline
\end{enumerate}
\end{hookbox}

\vspace{0.1in}

% ── DIFFERENCE IN PROPORTIONS ────────────────────────────────────────────────
\begin{notesbox}{Step 1 --- Difference in Proportions (UNC-1.Q, UNC-1.R)}
\small

\noindent\textbf{Formula:}
\[
  \hat{p}_1 - \hat{p}_2 = \frac{\text{successes in Group 1}}{\text{Group 1 total}} -
                           \frac{\text{successes in Group 2}}{\text{Group 2 total}}
\]

\vspace{4pt}
\textbf{Vaccine example} (Group 1 = Vaccinated, Group 2 = Unvaccinated):
\[
  \hat{p}_V - \hat{p}_U = \frac{\blank{1cm}}{\blank{1cm}} - \frac{\blank{1cm}}{\blank{1cm}}
  = \blank{2cm} - \blank{2cm} = \blank{2.4cm}
\]

\vspace{6pt}
\textbf{Interpretation template:}\\
\textit{``The proportion of [outcome] among [Group 1] was \underline{\hspace{2cm}} [higher / lower]
than the proportion among [Group 2].''}

\vspace{6pt}
\textbf{Write the interpretation for the vaccine example:}\\[6pt]
\writeline\\[1pt]\writeline

\vspace{6pt}
\textbf{Key properties:}
\begin{itemize}[leftmargin=1.5em, itemsep=2pt]
  \item Value of \blank{1cm} $\to$ no difference (no association).
  \item Positive $\to$ Group 1 rate is \blank{3cm}; Negative $\to$ Group 1 rate is \blank{3cm}.
  \item Units: \blank{3cm} (same units as the proportions themselves).
\end{itemize}
\end{notesbox}

\vspace{0.1in}

% ── RELATIVE RISK ─────────────────────────────────────────────────────────────
\begin{notesbox}{Step 2 --- Relative Risk (UNC-1.Q, UNC-1.R)}
\small

\noindent\textbf{Formula:}
\[
  \text{Relative Risk} = \frac{\hat{p}_1}{\hat{p}_2}
\]
The denominator $\hat{p}_2$ is the \textbf{\blank{3cm} group} (the comparison baseline).

\vspace{4pt}
\textbf{Vaccine example} (Group 1 = Vaccinated, Group 2 = Unvaccinated):
\[
  \text{RR} = \frac{\hat{p}_V}{\hat{p}_U} = \frac{\blank{2cm}}{\blank{2cm}} = \blank{2.4cm}
\]

\vspace{6pt}
\textbf{Interpretation template:}\\
\textit{``[Group 1] was \underline{\hspace{2cm}} times as likely to [outcome] as [Group 2].''}

\vspace{6pt}
\textbf{Write the interpretation for the vaccine example:}\\[6pt]
\writeline\\[1pt]\writeline

\vspace{6pt}
\textbf{Key properties:}
\begin{itemize}[leftmargin=1.5em, itemsep=2pt]
  \item Value of \blank{1cm} $\to$ no difference (no association).
  \item RR $> 1$ $\to$ Group 1 rate is \blank{3cm}; RR $< 1$ $\to$ Group 1 rate is \blank{3cm}.
  \item \textbf{Not} a percentage-point difference --- RR $= 2$ means ``twice as likely,'' not ``2\% more.''
\end{itemize}
\end{notesbox}

\vspace{0.1in}

% ── COMPARING THE TWO STATISTICS ─────────────────────────────────────────────
\begin{notesbox}{Step 3 --- When Do They Tell Different Stories?}
\small
\begin{multicols}{2}

\textbf{Context A: Rare outcome}\\
Group 1: 2 of 1000 ($\hat{p}_1 = 0.002$)\\
Group 2: 1 of 1000 ($\hat{p}_2 = 0.001$)

\vspace{4pt}
Difference in proportions: \blank{3cm}\\
Relative risk: \blank{3cm}

\vspace{6pt}
Which statistic is more striking?\\
\writeline\\[1pt]\writeline

\columnbreak

\textbf{Context B: Common outcome}\\
Group 1: 600 of 1000 ($\hat{p}_1 = 0.60$)\\
Group 2: 400 of 1000 ($\hat{p}_2 = 0.40$)

\vspace{4pt}
Difference in proportions: \blank{3cm}\\
Relative risk: \blank{3cm}

\vspace{6pt}
Which statistic is more striking?\\
\writeline\\[1pt]\writeline

\end{multicols}

\vspace{4pt}
\textbf{Takeaway:} Fill in the blanks.\\
A small \blank{3cm} can correspond to a large \blank{3cm}, and vice versa.
Always report \blank{5cm} in context.
\end{notesbox}

\vspace{0.1in}

% ── CONNECTIONS ───────────────────────────────────────────────────────────────
\begin{spiralbox}
\small
\begin{multicols}{2}

\textbf{How this connects:}
\begin{itemize}[itemsep=3pt]
  \item Lesson 2.2: conditional frequencies show \emph{whether} groups differ.
  \item Lesson 2.3 (today): difference in proportions and relative risk \emph{quantify} how much.
  \item Unit 6: formal inference tests whether the observed difference is statistically significant.
\end{itemize}

\columnbreak

\textbf{In your own words:}\\
Why do researchers sometimes prefer relative risk over the difference in proportions?\\[2pt]
\writeline\\[1pt]\writeline\\[1pt]\writeline

\end{multicols}
\end{spiralbox}

\vspace{0.1in}

\begin{notesbox}{Additional Notes}
\vspace{0pt}
\writeline\\\writeline\\\writeline\\\writeline\\\writeline\\\writeline
\end{notesbox}

\end{document}
